% ch05-market.tex – Chapter 5: market research and existing products (final round, Writer).
\ChapterMission{המשימה: למפות את המוצרים הקיימים ולהראות מה כבר אפשר לקנות ומה עדיין חסר}
\ChapterIcon{sensor}
\chapter{חקר שוק ומוצרים קיימים}\label{chap:market}
\begin{scenebox}
על עמוד הרמזור תלויים כבר שלושה דברים: ראש אות עם ספרות כתומות, לחצן מתכת שחוק וקופסה אפורה קטנה שמביטה מטה אל המדרכה.
כל אחד מהם הותקן בשנה אחרת ובידי ספק אחר. אף אחד מהם אינו יודע מה עושים האחרים.
\end{scenebox}

העמוד העמוס בתמונה מסכם את מצב השוק: הרכיבים קיימים, אך כל אחד פועל לבדו.
הפרק בוחן את המוצרים הקיימים בשלוש שכבות, משווה ביניהם על פי נתונים טכניים שפורסמו, ומזהה את המקום שבו הפרויקט משתלב.
נתוני היקף שוק או מחירים אינם מוצגים, משום שלא נמצא להם מקור שניתן לאמת.

\section{שלוש שכבות של שוק}
\paragraph{שכבת התצוגה.}
ראשי אות להולכי רגל עם ספירה לאחור הם מוצר מדף. התקן האמריקני מחייב אותם כשמופע הפינוי ארוך משבע שניות, ומגדיר את צבע הספרות ואת גודלן \cite{fhwa2025mutcd}. ראש מסחרי טיפוסי משלב יד, דמות הולך וספירה בת שתי ספרות, ומכיל בקרת התנגשויות פנימית שאינה מאפשרת להדליק יחד את האות להליכה ואת האות לעצירה \cite{current2023gtx}.

\paragraph{שכבת החישה.}
גלאי הולכי רגל בשפת המדרכה הם מחלקת מוצרים מוכרת, שרשויות תחבורה רוכשות לפי מפרט: לונדון, למשל, דורשת זיהוי של ממתין תוך חצי שנייה ועבודה גם בתאורה חלשה, ומשתמשת בגלאי כדי לבטל בקשה למופע הליכה כשהממתין כבר הלך \cite{tfl2015kerbside}.
סקירה של המעבדה הבריטית לחקר תחבורה (\en{TRL}) משווה את הטכנולוגיות: אינפרה-אדום, מיקרוגל, עיבוד תמונה, מזרני לחץ ומכ״ם \cite{trl2006ppr087}. בשטח, גלאים פסיביים צמצמו את מספר היוצאים לכביש באדום, אך דרשו כיול של אזור הזיהוי \cite{hughes2000}.

\paragraph{שכבת הבקרה והנתונים.}
כאן השוק כמעט ריק. הגלאים הקיימים משמשים לבקשת מופע הליכה, לביטולו או להארכת הפינוי להולכים איטיים \cite{hughes2000,tfl2015kerbside}. ספירת ממתינים ומדידת זמן המתנתם הודגמו באב-טיפוס מחקרי בלבד \cite{wickramasinghe2020}, ולא נמצא מוצר שמקצר את האדום לפי מספר הממתינים וזמן המתנתם.

\section{השוואה טכנית}
טבלה~\ref{tab:market} משווה את מחלקות המוצרים לפי מה שהן מספקות לפרויקט. ההשוואה מראה שתי מגמות: הגלאים המבוססים על מצלמה הם היחידים שיכולים לספור, אך הם מביאים איתם סוגיית פרטיות ורגישות לתאורה; והגלאים הפשוטים אמינים וזולים, אך מזהים נוכחות בלבד.

\begin{table}[ht]
  \centering
  \caption{השוואה טכנית של מחלקות מוצרים קיימות}
  \label{tab:market}
  \small
  \begin{tabularx}{\linewidth}{lXX}
    \toprule
    \headrow \thead{מחלקה} & \thead{נתונים שפורסמו} & \thead{משמעות לפרויקט} \\
    \midrule
    ראש ספירה בפינוי & ספרות בגובה \unit{9}{in} לפחות; הספק הספירה \unit{8}{W}; בקרת התנגשויות פנימית \cite{current2023gtx} & מתאים לספירה בפינוי; התזמון חייב להגיע מן הבקר \\
    גלאי מיקרוגל ואינפרה-אדום & מוצר ותיק וזול; קושי עם הולכים איטיים מאוד ועם יישור \cite{trl2006ppr087} & מזהה נוכחות בלבד, בלי ספירה \\
    מזרן לחץ & אזור זיהוי מוגדר היטב, אך הולכים אינם תמיד עומדים עליו \cite{trl2006ppr087} & לא מתאים לספירה ולמעבר עמוס \\
    גלאי ראייה סטריאו & אזור \unit{10\,$\times$\,3}{m}; הספק \unit{8.0}{W}; תאורה נוספת מתחת ל-\unit{15}{lux}; הזרמת וידאו \cite{agd2024agd645} & מועמד לחישה; מחייב עיבוד בקצה ובלי שמירת תמונות \\
    ראייה ממוחשבת לספירה & ספירה נכונה ב-\pct{72}; מדידת זמן המתנה \cite{wickramasinghe2020} & מוכיח היתכנות; דיוק עדיין נמוך מן הנדרש \\
    \bottomrule
  \end{tabularx}
\end{table}

\begin{projectbox}
הפרויקט אינו מפתח חומרה חדשה. הוא משתמש בגלאי מדף ובראש ספירה תקני, ומוסיף את מה שחסר בשוק: לוגיקה בבקר שמקצרת את האדום לפי מספר הממתינים וזמן המתנתם, בגבולות בטוחים, ושכבת נתונים שמעבירה ספירות בלבד.
\end{projectbox}
